%%%%%%%%%%%%%%%%%%%%%%%%%%%%%%%%%%%%%%%%%
% Medium Length Professional CV
% LaTeX Template
% Version 3.0 (December 17, 2022)
%
% This template originates from:
% https://www.LaTeXTemplates.com
%
% Author:
% Vel (vel@latextemplates.com)
%
% Original author:
% Trey Hunner (http://www.treyhunner.com/)
%
% License:
% CC BY-NC-SA 4.0 (https://creativecommons.org/licenses/by-nc-sa/4.0/)
%
%%%%%%%%%%%%%%%%%%%%%%%%%%%%%%%%%%%%%%%%%

%----------------------------------------------------------------------------------------
%	PACKAGES AND OTHER DOCUMENT CONFIGURATIONS
%----------------------------------------------------------------------------------------

\documentclass[
	%a4paper, % Uncomment for A4 paper size (default is US letter)
	11pt, % Default font size, can use 10pt, 11pt or 12pt
]{resume} % Use the resume class

\usepackage{ebgaramond} % Use the EB Garamond font
\usepackage[dvipdfm,linkcolor=blue,anchorcolor=blue,urlcolor=blue,colorlinks=true]{hyperref}
\urlstyle{same}
\usepackage{pxjahyper}

\name{Kan Kusakabe} % Your name to appear at the top

\address{Hokkaido University, Japan}
\address{kusakabe.kan.v5@elms.hokudai.ac.jp}
\address{\url{https://kankusakabe.github.io/Portfolio}}


\begin{document}

\begin{rSection}{Research Topic}
Human-Computer Interaction (HCI); Interaction Technique for Mobile Device; Gesture Design; Sensing Technique; Machine Learning (ML); 3D Printing; UI/UX.
\end{rSection}
\begin{rSection}{Fellow}
1. Hokkaido University EXEX Doctoral Fellowship. 2024--2027. \$3,000 annual and stipend.

2. HOKUDAI TECH GARAGE Project. 2 months (March-April 2022), \$ 2,000 per project.
\end{rSection}

%%%%%%%%%%%%	EDUCATION SECTION
\begin{rSection}{Education}
	
	\textbf{Hokkaido University, Japan} \hfill \textit{April 2024 ---} \\ 
	Doctor of Computer Science, Advisor: Daisuke Sakamoto 
	% Overall GPA: 5.678
    
    \textbf{Hokkaido University, Japan} \hfill \textit{April 2022 -- March 2024} \\ 
	Master of Computer Science, Advisor: Daisuke Sakamoto 
	% Overall GPA: 5.678

    \textbf{Hokkaido University, Japan} \hfill \textit{April 2020 -- March 2022} \\ 
	Bachelor of Computer Science, Advisor: Daisuke Sakamoto 
	% Overall GPA: 5.678

    \textbf{National Institute of Technology(KOSEN), Gifu College, Japan} \hfill \textit{April 2015 -- March 2020} \\ 
	Associate Degree of Engineering, Advisor: Deguchi Toshinori
	% Overall GPA: 5.678
    	
\end{rSection}

%%%%%%%%%%%%	WORK EXPERIENCE SECTION
\begin{rSection}{Papers}
    
	\begin{rSubsection}{Peer-reviewed International Conferences}{}{}{}
        \item Yuki Abe*, \underline{\textbf{Kan Kusakabe*}}, Myungguen Choi*, Daisuke Sakamoto, Tetsuo Ono. *Joint first authors. Understanding Usability of VR Pointing Methods with a Handheld-style HMD for Onsite Exhibitions. ACM CHI Conference on Human Factors in Computing Systems (CHI 2025), ACM, April 2025. (Conference acceptance rate: 27 \%; \textbf{Honorable Mentions: top 5 \% of accepted papers}). [\href{https://doi.org/10.1145/3706598.3713874}{doi}]
        
    \end{rSubsection}
    \begin{rSubsection}{Peer-reviewed Domestic Journals}{}{}{}
		\item \underline{\textbf{Kan Kusakabe}}, Daisuke Sakamoto, Tetsuo Ono. RingSense: Exploring User-defined Gestures for Phone Ring Holders. Transactions of the Information Processing Society of Japan, Vol. 66, No. 2, Feb. 2025. [\href{https://doi.org/10.20729/0002000028}{doi}]
    \end{rSubsection}
    \begin{rSubsection}{Peer-reviewed Domestic Conferences}{}{}{}
		\item \underline{\textbf{Kan Kusakabe}}, Daisuke Sakamoto, Tetsuo Ono. スマートフォン背面のジェスチャ入力を実現するスマホリング型デバイスの設計と実装. 第30回インタラクティブシステムとソフトウェアに関するワークショップ（WISS 2022），日本ソフトウェア科学会, pp.8-14，December 2022. (\textbf{Long-press, acceptance rate: 12.5\%}). [\href{https://www.wiss.org/WISS2022Proceedings/data/02.pdf}{link}]
	\end{rSubsection}
    
	\begin{rSubsection}{Other Publications (Non-peer-reviewed)}{}{}{}
		\item \underline{\textbf{Kan Kusakabe*}}, Yuki Abe*, Daisuke Sakamoto, Tetsuo Ono. *Joint first author. Game-2-X: 種類が異なるゲームプレイ間を繋ぐシステムの提案. 第31回インタラクティブシステムとソフトウェアに関するワークショップ（WISS 2023），日本ソフトウェア科学会, デモンストレーション発表. December 2023. [\href{https://www.wiss.org/WISS2023Proceedings/data/1-A03.pdf}{link}]
		\item \underline{\textbf{Kan Kusakabe}}, Myungguen Choi, Daisuke Sakamoto, Tetsuo Ono. 無段階調整インタフェースのためのハンドジェスチャによる操作手法の探索的研究. 情報処理学会 研究報告ヒューマンコンピュータインタラクション（HCI）, 2022-HCI-197(25), 1-8, March 2022. [\href{http://id.nii.ac.jp/1001/00217337/}{link}]
	\end{rSubsection}

    \begin{rSubsection}{Invited Talks}{}{}{}
		\item Public Lecture / Yokohama City MICE Next-Generation Development Initiative: CHI 2025 Symposium. April 26 2025. [\href{https://sigchi.jp/symposium/chi2025_symposium/}{url}]
	\end{rSubsection}

\end{rSection}

%%%%%%%%%%%%	TECHNICAL STRENGTHS SECTION
\begin{rSection}{Technical Strengths}
    \textbf{Domain}: Rapid Prototyping, 3D Print, Mobile Application, Machine-Learning (ML)

    \textbf{Programming Languages}:
    Python, C, C\#, JavaScript, PHP, Swift, Objective-C, Kotlin, Java, Arduino, Processing, Dart, LATEX

    \textbf{Frameworks \& Libraries}:
    Tensorflow, Pytorch, Next.js, Flutter, React, Unity
    
    \textbf{Platforms \& Development Environments}:
    Android, iOS, Web (Frontend, Backend, Database), Mac, Arduino Platform, Raspberry Pi

    \textbf{Design}:
    Adobe (Illustrator, Photoshop, Premiere Pro, After Effects, Audition, Figma), Fusion 360
\end{rSection}






\end{document}
